% Shared preamble for Calc 3 notes.
\usepackage[a4paper,margin=0.9in]{geometry}
\usepackage[utf8]{inputenc}
\usepackage[T1]{fontenc}
\usepackage{lmodern}
\usepackage{microtype}
\usepackage{amsmath,amssymb,amsthm,mathtools}
\usepackage{enumitem}
\usepackage{xcolor}
\usepackage[most]{tcolorbox}
\usepackage{hyperref}
\usepackage{titlesec}
\usepackage{parskip}
\usepackage{bm}

\hypersetup{
  colorlinks=true,
  linkcolor=black!70,
  urlcolor=blue!60!black,
  pdftitle={Calculus 3 -- Notes},
  pdfauthor={mathcheatsheet}
}

\theoremstyle{definition}
\newtheorem{definition}{Definition}[section]
\newtheorem{example}[definition]{Example}

\theoremstyle{plain}
\newtheorem{theorem}[definition]{Theorem}
\newtheorem{proposition}[definition]{Proposition}
\newtheorem{lemma}[definition]{Lemma}
\newtheorem{corollary}[definition]{Corollary}

\theoremstyle{remark}
\newtheorem*{remark}{Remark}
\newtheorem*{intuition}{Intuition}

\newtcolorbox{keybox}[1][]{
  enhanced, breakable, colback=blue!4, colframe=blue!55!black,
  fonttitle=\bfseries, title=#1, left=6pt, right=6pt, top=4pt, bottom=4pt,
  arc=2pt, boxrule=0.6pt
}
\newtcolorbox{pitfallbox}{
  enhanced, breakable, colback=red!4, colframe=red!55!black,
  fonttitle=\bfseries, title=Pitfalls, left=6pt, right=6pt, top=4pt, bottom=4pt,
  arc=2pt, boxrule=0.6pt
}

\newcommand{\R}{\mathbb{R}}
\newcommand{\N}{\mathbb{N}}
\newcommand{\eps}{\varepsilon}
\newcommand{\dx}{\,dx}
\newcommand{\dy}{\,dy}
\newcommand{\dz}{\,dz}
\newcommand{\dt}{\,dt}
\newcommand{\dA}{\,dA}
\newcommand{\dV}{\,dV}
\newcommand{\dS}{\,dS}
\newcommand{\dr}{\,dr}
\newcommand{\du}{\,du}
\newcommand{\dv}{\,dv}
\newcommand{\dth}{\,d\theta}
\newcommand{\dphi}{\,d\varphi}
\newcommand{\drho}{\,d\rho}
\newcommand{\vc}[1]{\mathbf{#1}}     % vector
\newcommand{\F}{\vc{F}}
\newcommand{\rr}{\vc{r}}
\newcommand{\nn}{\vc{n}}
\newcommand{\TT}{\vc{T}}
\newcommand{\NN}{\vc{N}}
\newcommand{\ii}{\vc{i}}
\newcommand{\jj}{\vc{j}}
\newcommand{\kk}{\vc{k}}
\newcommand{\dotp}{\boldsymbol\cdot}
\newcommand{\grad}{\nabla}
\DeclareMathOperator{\divg}{div}
\DeclareMathOperator{\curl}{curl}
\DeclareMathOperator{\proj}{proj}

\titleformat{\section}{\Large\bfseries\color{blue!40!black}}{\thesection}{0.6em}{}
\titleformat{\subsection}{\large\bfseries\color{blue!25!black}}{\thesubsection}{0.5em}{}
